% GENERATED FILE. Edit canonical lessons, then run npm run generate:books.
% canonical-source-hash: fnv1a64:e0c8b69641c8b22f
% canonical-lessons: GE-C185-retten, GE-C185-rudern, GE-C185-schieben, GE-C185-schlagen, GE-C185-schmecken

\chapter{To rescue, To row, To push, To hit, To taste}
\label{ch:ge-to-rescue-to-row-to-push-to-hit-to-taste}
\hlchaptermodality{185}

\begin{chapteropening}
\textbf{By the end of this chapter:} \emph{I can say to rescue, to row, to push, to hit and to taste.}

Complete the last lesson of chapter 185: I can say to rescue, to row, to push, to hit and to taste.
\end{chapteropening}

\section[retten]{retten — to rescue}

\label{lesson:GE-C185-retten}



\begin{quote}
Before the new one: say the German for to test, then the German for to clean.
\end{quote}

\subsection*{You'll want to know: retten}
\textbf{retten} — \textquotedblleft{}to rescue\textquotedblright{}.

\begin{grammarlens}[title={What you've built}]
One more word for this chapter.
\end{grammarlens}

\subsection*{Your turn}
\begin{itemize}
\raggedright
  \item \emph{Say it:} \emph{retten}
  \item \emph{Say it:} \emph{retten}, once more
  \item \emph{Say it:} say \emph{retten}
  \item \emph{Recall:} say the German for to test, then the German for to clean, then say \emph{retten} again
\end{itemize}

\begin{tcolorbox}[breakable,colback=teal!4,colframe=teal!35!black,title={Before you move on}]
What does retten mean? (\textquotedblleft{}To rescue\textquotedblright{}.) Say it once more.
\end{tcolorbox}
\section[rudern]{rudern — to row}

\label{lesson:GE-C185-rudern}



\begin{quote}
Before the new one: say the German for to clean, then the German for to rescue.
\end{quote}

\subsection*{You'll want to know: rudern}
\textbf{rudern} — \textquotedblleft{}to row\textquotedblright{}.

\begin{grammarlens}[title={What you've built}]
One more word for this chapter.
\end{grammarlens}

\subsection*{Your turn}
\begin{itemize}
\raggedright
  \item \emph{Say it:} \emph{rudern}
  \item \emph{Say it:} \emph{rudern}, once more
  \item \emph{Say it:} say \emph{rudern}
  \item \emph{Recall:} say the German for to clean, then the German for to rescue, then say \emph{rudern} again
\end{itemize}

\begin{tcolorbox}[breakable,colback=teal!4,colframe=teal!35!black,title={Before you move on}]
What does rudern mean? (\textquotedblleft{}To row\textquotedblright{}.) Say it once more.
\end{tcolorbox}
\section[schieben]{schieben — to push}

\label{lesson:GE-C185-schieben}



\begin{quote}
Before the new one: say the German for to rescue, then the German for to row.
\end{quote}

\subsection*{You'll want to know: schieben}
\textbf{schieben} — \textquotedblleft{}to push\textquotedblright{}.

\begin{grammarlens}[title={What you've built}]
One more word for this chapter.
\end{grammarlens}

\subsection*{Your turn}
\begin{itemize}
\raggedright
  \item \emph{Say it:} \emph{schieben}
  \item \emph{Say it:} \emph{schieben}, once more
  \item \emph{Say it:} say \emph{schieben}
  \item \emph{Recall:} say the German for to rescue, then the German for to row, then say \emph{schieben} again
\end{itemize}

\begin{tcolorbox}[breakable,colback=teal!4,colframe=teal!35!black,title={Before you move on}]
What does schieben mean? (\textquotedblleft{}To push\textquotedblright{}.) Say it once more.
\end{tcolorbox}
\section[schlagen]{schlagen — to hit, to beat}

\label{lesson:GE-C185-schlagen}



\begin{quote}
Before the new one: say the German for to row, then the German for to push.
\end{quote}

\subsection*{You'll want to know: schlagen}
\textbf{schlagen} — \textquotedblleft{}to hit, to beat\textquotedblright{}.

\begin{grammarlens}[title={What you've built}]
One more word for this chapter.
\end{grammarlens}

\subsection*{Your turn}
\begin{itemize}
\raggedright
  \item \emph{Say it:} \emph{schlagen}
  \item \emph{Say it:} \emph{schlagen}, once more
  \item \emph{Say it:} say \emph{schlagen}
  \item \emph{Recall:} say the German for to row, then the German for to push, then say \emph{schlagen} again
\end{itemize}

\begin{tcolorbox}[breakable,colback=teal!4,colframe=teal!35!black,title={Before you move on}]
What does schlagen mean? (\textquotedblleft{}To hit, to beat\textquotedblright{}.) Say it once more.
\end{tcolorbox}
\section[schmecken]{schmecken — to taste}

\label{lesson:GE-C185-schmecken}



\begin{quote}
Before the new one: say the German for to push, then the German for to hit.
\end{quote}

\subsection*{You'll want to know: schmecken}
\textbf{schmecken} — \textquotedblleft{}to taste\textquotedblright{}.

\begin{grammarlens}[title={What you've built}]
One more word for this chapter.
\end{grammarlens}

\subsection*{Your turn}
\begin{itemize}
\raggedright
  \item \emph{Say it:} \emph{schmecken}
  \item \emph{Say it:} \emph{schmecken}, once more
  \item \emph{Say it:} say \emph{schmecken}
  \item \emph{Recall:} say the German for to push, then the German for to hit, then say \emph{schmecken} again
\end{itemize}

\begin{tcolorbox}[breakable,colback=teal!4,colframe=teal!35!black,title={Before you move on}]
What does schmecken mean? (\textquotedblleft{}To taste\textquotedblright{}.) Say it once more.
\end{tcolorbox}
